\subsubsection{\texorpdfstring{\(\Gamma\)}{Gamma} 重整化与端点层尾项账本}\label{subsubsec:typed-address-biaxial-completion-gamma-renormalized-endpoint-budget}
固定原点 Taylor 证书的端点成本由两部分叠加而成：一部分来自最近解析障碍 \(w=1\) 的几何距离，另一部分来自 \(F_a=\xi\circ s_a\) 中 \(\Gamma(s/2)\) 的快速增长。后者不携带零点信息，因而应先从证书对象中剥离。定义
\[
\widetilde F_a(w)
:=
F_a(w)\frac{\pi^{s_a(w)/2}}{\Gamma(s_a(w)/2)}
=
\frac12s_a(w)(s_a(w)-1)\zeta(s_a(w)).
\]
由于 \(\pi^{s/2}/\Gamma(s/2)\) 在 \(\Re s>1/2\) 内解析且无零点，\(\widetilde F_a\) 与 \(F_a\) 在单位圆盘内具有完全相同的零点；其中 \(s=1\) 处的 \(\zeta\)-极点由因子 \(s-1\) 抵消，给出可去奇点。

\begin{lemma}[Cayley 图表的端点几何界]\label{lem:typed-address-biaxial-completion-cayley-endpoint-geometry-bound}
固定 \(a>0\)、\(0<\eta<1\)，并令 \(0<\varepsilon\le1/2\)。若
\[
|w|\le 1-\eta\varepsilon,\qquad
u(w):=\frac{1+w}{1-w},\qquad
s_a(w)=\frac12+au(w),
\]
则
\[
\Re u(w)=\frac{1-|w|^2}{|1-w|^2}>0,
\]
且存在只依赖于 \(a,\eta\) 的常数 \(C_1,C_2>0\)，使
\[
|s_a(w)|\le C_1\varepsilon^{-1},
\qquad
\Re s_a(w)\ge \frac12+C_2\varepsilon .
\]
\end{lemma}
\begin{proof}
令 \(\rho=1-\eta\varepsilon\)。三角不等式给出
\[
|u(w)|\le \frac{1+\rho}{1-\rho}
=\frac{2-\eta\varepsilon}{\eta\varepsilon}
\le \frac{2}{\eta\varepsilon}.
\]
另一方面，
\[
\Re u(w)=\frac{1-|w|^2}{|1-w|^2}
\ge
\frac{1-\rho^2}{(1+\rho)^2}
=
\frac{1-\rho}{1+\rho}
=
\frac{\eta\varepsilon}{2-\eta\varepsilon}.
\]
代入 \(s_a(w)=1/2+au(w)\) 即得所示两个界。证毕。
\end{proof}

\begin{proposition}[重整化函数的端点多项式增长]\label{prop:typed-address-biaxial-completion-renormalized-polynomial-endpoint-growth}
固定 \(a>0\) 与 \(\eta\in(0,1)\)。存在常数 \(C_{a,\eta}>0\) 与 \(B>0\)，使得对所有 \(0<\varepsilon\le1/2\) 有
\[
\sup_{|w|\le 1-\eta\varepsilon}
|\widetilde F_a(w)|
\le
C_{a,\eta}\varepsilon^{-B}.
\]
\end{proposition}
\begin{proof}
记 \(s=s_a(w)\)。由引理 \ref{lem:typed-address-biaxial-completion-cayley-endpoint-geometry-bound}，
\[
|s|\le C_1\varepsilon^{-1},
\qquad
\Re s>\frac12.
\]
若 \(\Re s\ge2\)，则 \(|\zeta(s)|\le\zeta(2)\)，故
\[
|\widetilde F_a(w)|
\le
\frac12\zeta(2)|s(s-1)|
\le
C_{a,\eta}\varepsilon^{-2}.
\]
若 \(1/2<\Re s\le2\)，则 \(H(s):=(s-1)\zeta(s)\) 在闭竖条 \(1/2\le\Re s\le2\) 中全纯，并由 \(\zeta\)-函数的标准竖条增长估计满足
\[
|H(s)|\le C(1+|\Im s|)^A
\]
对某些绝对常数 \(A,C>0\) 成立；这可由经典 Euler--Maclaurin 展开、函数方程与 Phragm\'en--Lindel\"of 原理推出，亦见 \cite{Titchmarsh1986RiemannZeta}。结合 \(|\Im s|\le |s|\le C_1\varepsilon^{-1}\)，得到
\[
|\widetilde F_a(w)|
=
\frac12|sH(s)|
\le
C_{a,\eta}\varepsilon^{-(A+1)}.
\]
取 \(B=\max\{2,A+1\}\)，并调整常数，得到所示一致上界。证毕。
\end{proof}

令
\[
\widetilde G_{a,q}(z):=\widetilde F_a(r_qz),
\qquad
r_q:=1-2^{-q},
\qquad
\varepsilon_q:=2^{-q},
\]
并取
\[
R_{q,\eta}:=\frac{1-\eta\varepsilon_q}{r_q}.
\]
则 \(r_qR_{q,\eta}=1-\eta\varepsilon_q<1\)，且
\[
R_{q,\eta}-1
=
\frac{(1-\eta)\varepsilon_q}{1-\varepsilon_q}
\asymp_{\eta}\varepsilon_q,
\qquad
\log R_{q,\eta}\asymp_{\eta}\varepsilon_q.
\]

\begin{proposition}[重整化 Taylor 尾项账本]\label{prop:typed-address-biaxial-completion-renormalized-taylor-tail-budget}
设
\[
\widetilde P_{q,N}(z)=\sum_{k=0}^{N}\widetilde c_k^{(q)}z^k
\]
为 \(\widetilde G_{a,q}\) 的 \(N\) 阶 Taylor 截断，并记
\[
\widetilde E_{q,N}:=
\sup_{|z|=1}|\widetilde G_{a,q}(z)-\widetilde P_{q,N}(z)|,
\qquad
\widetilde m_q:=\min_{|z|=1}|\widetilde G_{a,q}(z)|.
\]
若
\[
M_{q,\eta}:=
\sup_{|z|=R_{q,\eta}}|\widetilde G_{a,q}(z)|,
\]
则
\[
\widetilde E_{q,N}
\le
\frac{M_{q,\eta}R_{q,\eta}^{-N}}{R_{q,\eta}-1},
\qquad
M_{q,\eta}\le C_{a,\eta}2^{Bq}.
\]
因此，若给定目标精度 \(\mu>0\)，则 \(\widetilde E_{q,N}\le\mu\) 由充分条件
\[
N>
\frac{
\log C_{a,\eta}
+(B+1)q\log2
+\log(1/\mu)
}{
\log R_{q,\eta}
}
\]
保证。特别地，取 \(\mu=\widetilde m_q/2\)，便得到一个足以闭合 Rouch\'e 误差的阶数条件
\[
N
\ge
C'_{a,\eta}2^q
\left(
\log\frac1{\widetilde m_q}
+(B+1)q\log2
+1
\right).
\]
若 \(\log(1/\widetilde m_q)=O(q)\)，则可取 \(N=O_{a,\eta}(q2^q)\)。
\end{proposition}
\begin{proof}
Cauchy 尾项估计给出第一式。由 \(r_qR_{q,\eta}=1-\eta\varepsilon_q\) 与命题 \ref{prop:typed-address-biaxial-completion-renormalized-polynomial-endpoint-growth}，
\[
M_{q,\eta}\le C_{a,\eta}\varepsilon_q^{-B}=C_{a,\eta}2^{Bq}.
\]
又有 \(R_{q,\eta}-1\asymp_{\eta}2^{-q}\) 与 \(\log R_{q,\eta}\asymp_{\eta}2^{-q}\)。将这些界代入尾项估计，得
\[
\widetilde E_{q,N}
\le
C_{a,\eta}2^{(B+1)q}R_{q,\eta}^{-N}.
\]
要求右端不超过 \(\mu\)，取对数即得第一个阶数条件。再令 \(\mu=\widetilde m_q/2\)，并用 \((\log R_{q,\eta})^{-1}\le C'_\eta2^q\)，得到所示 dyadic 条件。证毕。
\end{proof}

等价地，若写 \(r=1-\varepsilon\) 而非 dyadic 半径，则同一估计为
\[
\widetilde E_{r,N}
\le
C_{a,\eta}\varepsilon^{-(B+1)}R_{\varepsilon,\eta}^{-N},
\qquad
R_{\varepsilon,\eta}:=\frac{1-\eta\varepsilon}{1-\varepsilon}.
\]
因此
\[
N
\ge
C_{a,\eta}\varepsilon^{-1}
\left(
\log\frac1{\widetilde m_r}
+(B+1)\log\frac1\varepsilon
+1
\right)
\]
足以保证 \(\widetilde E_{r,N}<\widetilde m_r\)。结合端点压缩律 \(\varepsilon\sim\gamma^{-2}\)，理想最小模制度下的固定原点证书尺度正是
\[
N\sim \gamma^2\log|\gamma|.
\]
由此可见，\(\Gamma\)-重整化已经去除了原始 \(\xi\)-函数中的超额增长成本；端点层中剩余的非几何瓶颈集中在边界最小模 \(\widetilde m_q\) 及其可审计下界。

重整化后的固定半径证书因而采用
\[
\widetilde{\mathsf{SC}}(q,N)
=
\bigl(\widetilde P_{q,N},\,\widetilde E_{q,N}^+,\,\widetilde m_q^-,\,\widetilde\lambda_{q,N}^-\bigr),
\]
其中 \(\widetilde\lambda_{q,N}^-\) 是
\(\mathfrak C(\widetilde P_{q,N})\) 的严格谱裕度下界。若
\[
\widetilde E_{q,N}^+<\widetilde m_q^-,
\qquad
\widetilde\lambda_{q,N}^->0,
\]
则 \(\widetilde F_a\) 在 \(|w|<r_q\) 内无零点，进而 \(F_a\) 在同一子盘内无零点。该账本把端点层成本分离为三个量：几何分辨率 \(2^q\)、最小模惩罚 \(\log(1/\widetilde m_q)\)，以及已被多项式控制吸收的重整化增长项。
